% A geometric synopsis, indexed by section range.
% Nuclear colors follow the manuscript's chemistry palette.
\begin{figure}[!ht]
\centering
\begin{tikzpicture}[x=1cm,y=1cm,font=\small,
  line cap=round,line join=round,
  stageheading/.style={font=\footnotesize\bfseries,anchor=west},
  note/.style={font=\scriptsize,text=figMuted,align=center},
  flow/.style={-{Stealth[length=2mm]},draw=figMuted,line width=.7pt},
  pics/nucleiset/.style={code={
    \foreach \x/\y in {-.85/-.30,-.85/.30,-.30/.55,.85/.30,.85/-.30}
      \draw[figInk,fill=figHydrogen,line width=.45pt]
        (\x,\y) circle[radius=.075cm];
    \fill[figCarbon] (-.30,0) circle[radius=.105cm];
    \fill[figNitrogen] (.30,0) circle[radius=.105cm];
  }}]
  \path[use as bounding box] (0,.2) rectangle (14.4,11.3);

  % 1. Configurations carry spin vectors before a molecular grouping is chosen.
  \node[stageheading] at (.1,10.95) {\S\S\,1--3\quad Nuclei and spin};
  \path[fill=figInk!3,rounded corners=12pt]
    (.1,7.30) rectangle (6.55,8.70);
  \node at (.45,8.43) {$\conf$};
  \pic[scale=.66,rotate=20] at (1.2,7.95) {nucleiset};
  \pic[scale=.66] at (3.35,7.95) {nucleiset};
  \pic[scale=.66,rotate=-40] at (5.5,7.95) {nucleiset};
  \node[note] at (1.2,7.48) {$X_1$};
  \node[font=\scriptsize] at (3.35,7.48) {$X$};
  \node[note] at (5.5,7.48) {$X_2$};
  \node[note,text=figBlue] at (3.35,10.36) {a copy of $\spin$ over each configuration};
  \foreach \x/\dx/\dy in {1.2/.10/.42,3.35/.43/.32,5.5/.55/.17} {
    \draw[figLine,densely dotted,line width=.65pt]
      (\x,8.43)--(\x,9.25);
    \path[draw=figBlue!60,fill=figBlue!5,line width=.55pt]
      (\x-.65,9.15)--(\x+.30,9.15)--(\x+.65,9.90)--(\x-.30,9.90)--cycle;
    \fill[figBlue] (\x-.20,9.35) circle[radius=1pt];
    \draw[figBlue,line width=.85pt,-{Stealth[length=1.5mm]}]
      (\x-.20,9.35)--++(\dx,\dy);
  }
  \node[font=\scriptsize,text=figBlue,anchor=south] at (3.73,9.80) {$\Psi(X)$};
  \node[font=\scriptsize] at (3.3,6.85)
    {$\Psi\circ\sigma\sim\stat(\sigma)\,\sigma\act\Psi$};

  % 2. Grouping is extra structure. A single fragment gives a molecule.
  \node[stageheading] at (7.7,10.95) {\S\S\,4--6\quad Groupings and symmetry};
  \draw[figTeal,fill=figTeal!5] (8.65,9.30)
    ellipse[x radius=.57cm,y radius=.71cm];
  \draw[figTeal,fill=figTeal!5] (10.02,9.30)
    ellipse[x radius=.51cm,y radius=.71cm];
  \fill[figCarbon] (8.80,9.3) circle[radius=.095cm];
  \foreach \x/\y in {8.35/9.05,8.35/9.55,8.80/9.85}
    \draw[figInk,fill=figHydrogen,line width=.45pt]
      (\x,\y) circle[radius=.07cm];
  \fill[figNitrogen] (9.88,9.30) circle[radius=.095cm];
  \foreach \y in {9.02,9.58}
    \draw[figInk,fill=figHydrogen,line width=.45pt]
      (10.25,\y) circle[radius=.07cm];
  \draw[flow] (10.75,9.3)--(11.30,9.3);
  \draw[figTeal,fill=figTeal!5] (12.65,9.35)
    ellipse[x radius=1.13cm,y radius=.75cm];
  \begin{scope}[shift={(12.65,9.3)}]
    \draw[figLine] (-.85,-.3)--(-.3,0)--(-.85,.3)
      (-.3,.55)--(-.3,0)--(.3,0)--(.85,.3)
      (.3,0)--(.85,-.3);
    \pic {nucleiset};
  \end{scope}
  \node[note] at (9.3,8.28) {several fragments};
  \node[note] at (12.65,8.28) {one molecule};
  \node (ah) at (8.35,7.32) {$H$};
  \node (ag6) at (10.48,7.32) {$G_6$};
  \node (ag12) at (12.65,7.32) {$G_{12}$};
  \draw[flow,figTeal] (ah)--node[above,font=\scriptsize] {$t$} (ag6);
  \draw[flow,figGold,dashed] (ag6)--node[above,font=\scriptsize] {$tu$} (ag12);
  \node[note] at (10.5,6.70) {resolve rearrangements};

  % 3. Soft localization and an optional exact support restriction.
  \node[stageheading] at (.1,5.45) {\S\S\,7--9\quad Localized states and geometry};
  \pic[scale=1.08] at (2.0,2.80) {versiontube};
  \foreach \angle/\height in {-30/0,90/0,210/0,150/1.3,270/1.3,30/1.3} {
    \coordinate (packet) at
      ({2+1.08*1.34*cos(\angle)},{2.8+1.08*(\height+.43*sin(\angle))});
    \draw[figTeal,opacity=.23,line width=.5pt]
      (packet) ellipse[x radius=.29cm,y radius=.17cm];
    \draw[figTeal,opacity=.15,line width=.5pt]
      (packet) ellipse[x radius=.43cm,y radius=.26cm];
  }
  \node[text=figTeal] at (2.0,1.68) {$\mathcal F$};
  \node[note] at (2.0,1.22) {one connected family};
  \node[note] at (2.0,.72) {versions $G/H$, copies $S/G$};
  \node[note] at (5.2,4.65) {soft localization};
  \draw[figLine] (4.18,3.55)--(6.3,3.55);
  \draw[figTeal,line width=.85pt]
    plot[domain=-1.05:1.05,samples=60] ({5.24+\x},{3.55+.65*exp(-5*\x*\x)});
  \draw[flow] (5.24,3.18)--(5.24,2.83);
  \node[note,anchor=west] at (5.43,3.00) {optional};
  \foreach \x in {4.54,5.94}
    \draw[figGold,dashed,line width=.6pt] (\x,1.73)--(\x,2.64);
  \draw[figLine] (4.18,1.87)--(6.3,1.87);
  \draw[figTeal,line width=.85pt]
    (4.18,1.87)--(4.54,1.87)
    plot[domain=-.99:.99,samples=80]
      ({5.24+.70*\x},{1.87+.65*exp(1-1/(1-\x*\x))})
    --(6.3,1.87);
  \node[note] at (5.24,1.22) {supported packets};

  % 4. Charts describe the connected family; the mode expansion changes basis.
  \node[stageheading] at (7.7,5.45) {\S\S\,10--12\quad Position and momentum};
  \path[draw=figLine,fill=figTeal!5,rounded corners=3pt]
    (7.90,3.10) rectangle (9.53,4.20);
  \path[draw=figLine,fill=figBlue!5,rounded corners=3pt]
    (8.72,3.40) rectangle (10.35,4.50);
  \draw[figTeal,line width=1pt,-{Stealth[length=1.5mm]}]
    (8.05,3.50)..controls(8.85,3.42)and(9.1,4.40)..(10.17,3.87);
  \node[note] at (9.15,4.86) {overlapping charts};
  \node[font=\scriptsize] at (9.12,2.55) {$|r,q,Q\rangle$};
  \draw[flow] (10.55,3.90)--(11.03,3.90);
  \foreach \level/\mode/\shade/\pattern in
    {3.14/1/figEtwo/dashed,3.64/2/figEtwo/dashed,4.14/3/figAtwo/densely dotted} {
    \draw[figLine,line width=.35pt] (11.3,\level)--(13.8,\level);
    \draw[\shade,\pattern,line width=.85pt]
      plot[domain=0:1,samples=65] ({11.3+2.5*\x},{\level+.18*sin(360*\mode*\x)});
  }
  \foreach \x in {11.3,13.8}
    \draw[figMuted,densely dashed,line width=.45pt] (\x,2.92)--(\x,4.40);
  \node[note] at (12.55,4.86) {adapted modes};
  \node[font=\scriptsize] at (12.55,2.55) {$|JMK,m,p\rangle$};
  \node[note] at (9.15,1.80) {gluing and monodromy};
  \node[note] at (12.55,1.80) {species and spin};
  \node[note] at (11,1.22) {the same states in a familiar basis};
\end{tikzpicture}
\caption*{\textbf{Geometric abstract.} Each configuration $X$ carries a spin
fiber $\spin$, with value $\Psi(X)$. Choose groupings and feasible motions.
Build localized states and feasible families, then describe the same states
in position and momentum. Hard support is optional. Chart gluing preserves
the distinction between separate feasible components.}
\end{figure}
